\section*{Contoh}
\addcontentsline{toc}{section}{Contoh}

\begin{contoh}[Membaca anatomi satu bab]
Perhatikan sebuah bab inti, misalnya Bab VII tentang implementasi
lexer. Kotak outcome menyebut Sub-CPMK 2.4, yaitu mengimplementasikan
lexical analyzer untuk bahasa sederhana. Bagian materi menjelaskan
buffer, sentinel, dan \eng{maximal munch}. Contoh menunjukkan potongan
kode \code{next\_token}. Praktikum meminta mahasiswa melengkapi lexer
dan menjalankan 30 tes. Asesmen menilai kebenaran, kualitas kode,
kelengkapan, dan pengujian. Daftar cek menyatakan bukti yang harus
ditunjukkan.

Dengan membaca anatomi ini, mahasiswa mengetahui sejak awal apa yang
harus dihasilkan, bukan hanya apa yang harus dibaca.
\end{contoh}

\begin{contoh}[Dari program sumber ke token]
Untuk program pada \lstref{lst:ch01-hello}, urutan token yang dihasilkan
lexer kira-kira adalah:
\begin{center}
\code{int}, \code{main}, \code{(}, \code{)}, \code{\{}, \code{int},
\code{a}, \code{=}, \code{5}, \code{;}, \code{int}, \code{b},
\code{=}, \code{7}, \code{;}, \code{print}, \code{(}, \code{a},
\code{+}, \code{b}, \code{)}, \code{;}, \code{return}, \code{0},
\code{;}, \code{\}}, \code{EOF}.
\end{center}
Perhatikan bahwa spasi dan komentar tidak menghasilkan token, sedangkan
setiap tanda baca dan literal menghasilkan token tersendiri.
\end{contoh}
